\protect \section *{Розділ~\ref *{sec:mhd-waves}. \nameref *{sec:mhd-waves}}
\begin{Solution}{6.{1}}
		$v_A\approx7{,}7\times10^{8}$~см/с (дейтерій), $c_s\approx1{,}4\times10^{8}$~см/с, $v_{Ti}\sim10^{8}$~см/с~--- $v_A\gg v_{Ti}$; час проходу $a/v_A\approx2{,}6\times10^{-7}$~с.
	
\end{Solution}
\begin{Solution}{6.{2}}
		при $\theta\to0$ $v_+\to\max(v_A,c_s)$, $v_-\to\min(v_A,c_s)$: для $v_A>c_s$ швидка мода
	переходить в~альвенівську, а~повільна~--- у~звук; для $v_A<c_s$~--- навпаки.
	
\end{Solution}
\begin{Solution}{6.{3}}
		$v_g=2c\sqrt{\omega}\,(\omega_{ce}-\omega)^{3/2}/(\omega_{pe}\omega_{ce})\approx3{,}2\times10^{9}$ і~$3{,}8\times10^{9}$~см/с;
		час проходу $\approx3{,}1$ і~$2{,}6$~с, $\Delta t\approx0{,}5$~с. Формула $v_g\approx2c\sqrt{\omega\omega_{ce}}/\omega_{pe}$ (границя $\omega\ll\omega_{ce}$) дала б
		$3{,}8\times10^{9}$ і~$5{,}4\times10^{9}$~см/с, $\Delta t\approx0{,}8$~с, тобто завищує $v_g$ приблизно на $17$ і~$40$~\%.
	
\end{Solution}
\begin{Solution}{6.{4}}
		$f_{pe}\approx63$~ГГц, $f_{ce}\approx84$~ГГц, $f_{UH}\approx105$~ГГц; $f_{ci}\approx46$~МГц,
	$f_{pi}\approx1{,}5$~ГГц, $f_{LH}\approx1{,}2$~ГГц. З~боку слабкого поля $X$-мода спершу натрапляє
	на відсікання $\omega=\omega_R$ і~відбивається; запущена з~боку сильного поля, вона доходить до резонансу.
	
\end{Solution}
\begin{Solution}{6.{5}}
		$\approx520$~км/с; $\approx220$~км/с; $\approx6{,}5$~км/с.
	
\end{Solution}
\begin{Solution}{6.{6}}
		у~короні ($T\sim100$~еВ) спітцерівська $\nu_m\sim10^{4}$~см$^2$/с; для $\omega\sim10^{-2}$~с$^{-1}$
	і~$v_A\sim7\times10^{7}$~см/с довжина загасання $1/\operatorname{Im}k\sim10^{24}$~см$\;\gg R_\odot$. Дисипація
	можлива лише на дуже малих масштабах (фазове перемішування, турбулентний каскад).
	
\end{Solution}
\begin{Solution}{6.{7}}
		$f_{pe}\approx90$~ГГц, $f_{ce}\lesssim140$~ГГц $\ll1{,}5$~ТГц; $\psi\approx0{,}53$~рад $\approx30^\circ$.
	
\end{Solution}
\begin{Solution}{6.{8}}
		$B=5$~Тл, $R=B_0R_0/B=3{,}6$~м.
	
\end{Solution}
\begin{Solution}{6.{9}}
		Вивід за~\cite{chen2016introduction}. Збурення беремо у~вигляді $\propto e^{i(kx-\omega t)}$, тож
		$\partial_t\to-i\omega$, $\partial_x\to ik$. Незбурена плазма квазінейтральна: $n_{e0}=Zn_{i0}$.

		\emph{Електрони.} Через $\omega\ll kv_{Te}$ вони больцманівські, $n_e=n_{e0}e^{e\varphi/T_e}$, і~після
		лінеаризації
		\begin{equation*}
			\frac{n_{e1}}{n_{e0}}=\frac{e\varphi}{T_e}.
		\end{equation*}

		\emph{Іони.} Сила на~іон $-Ze\,\partial_x\varphi$, а~тиск іонів $p_{i1}=\gamma_iT_in_{i1}$ з~$\gamma_i=3$, тож
		градієнт тиску на~один іон дорівнює $-\gamma_iT_i\,\partial_xn_{i1}/n_{i0}$:
		\begin{equation*}
			-i\omega m_iv=-ik\,Ze\varphi-ik\,3T_i\,\frac{n_{i1}}{n_{i0}}.
		\end{equation*}
		Неперервність $-i\omega n_{i1}+ikn_{i0}v=0$ дає $n_{i1}/n_{i0}=kv/\omega$, тобто $v=\omega n_{i1}/(kn_{i0})$.
		Підставляємо $v$ в~рівняння руху, множимо на~$ik$ і~групуємо члени з~$n_{i1}$:
		\begin{equation*}
			\Bigl(\omega^2-\frac{3k^2T_i}{m_i}\Bigr)\frac{n_{i1}}{n_{i0}}=\frac{k^2Ze\varphi}{m_i}
			\;\Rightarrow\;
			\frac{n_{i1}}{n_{i0}}=\frac{k^2Ze\varphi}{m_i\,(\omega^2-3k^2T_i/m_i)}.
		\end{equation*}

		\emph{Пуассон.} $\nabla^2\varphi=-4\pi\rho$ з~$\rho=e(Zn_{i1}-n_{e1})$ дає $k^2\varphi=4\pi e(Zn_{i1}-n_{e1})$.
		Підставляємо $n_{i1}$ та~$n_{e1}=Zn_{i0}\,e\varphi/T_e$ і~скорочуємо на~$\varphi$:
		\begin{equation*}
			k^2=4\pi Z^2e^2n_{i0}\left[\frac{k^2}{m_i(\omega^2-3k^2T_i/m_i)}-\frac1{ZT_e}\right].
		\end{equation*}
		Оскільки $4\pi Z^2e^2n_{i0}/(ZT_e)=4\pi n_{e0}e^2/T_e=\lambda_{De}^{-2}$, множимо на~$\lambda_{De}^2$:
		\begin{equation*}
			k^2\lambda_{De}^2+1=\frac{ZT_e}{m_i}\,\frac{k^2}{\omega^2-3k^2T_i/m_i}.
		\end{equation*}
		Звідси $\omega^2-3k^2T_i/m_i=k^2ZT_e/[m_i(1+k^2\lambda_{De}^2)]$, тобто
		\begin{equation*}
			\omega^2=k^2\Bigl[\frac{ZT_e}{m_i\,(1+k^2\lambda_{De}^2)}+\frac{3T_i}{m_i}\Bigr].
		\end{equation*}
	
\end{Solution}
\begin{Solution}{6.{10}}
		$c_s\approx3{,}1\times10^{6}$~см/с; $k\lambda_{De}\approx0{,}48$.
	
\end{Solution}
